% Overview and four visual steps for a network-security audience with basic FL/ML knowledge.
% Included by hivemind-control-plane.tex; shares its styles and speaker notes.

\begin{frame}{Inside a Hivemind peer}
\begin{visual}
  % An overview of components, with only the language boundary connected.
  \fill[ChalmersPurple!4,rounded corners=3mm]
    (1.8,2.1) rectangle (12.6,5.25);
  \node[anchor=west,font=\small\bfseries,text=ChalmersPurple]
    at (2.2,4.94) {Hivemind\quad\textbar\quad Python};

  \foreach \xx/\yy/\name/\purpose in {
    4.55/4.23/{All-reduce}/{Averages group gradients.},
    9.85/4.23/{Optimizer}/{Schedules averaging.},
    4.55/3.2/{Hivemind DHT}/{Stores peer metadata.},
    9.85/3.2/{Matchmaking}/{Forms temporary groups.}} {
    \node[box,draw=none,fill=white,minimum width=4.6cm,minimum height=0.85cm,
      inner sep=0pt] at (\xx,\yy) {};
    \node[font=\small\bfseries,text=ChalmersPurple] at (\xx,\yy+0.17) {\name};
    \node[quiet,text=ink] at (\xx,\yy-0.18) {\purpose};
  }

  \node[box,draw=none,fill=ChalmersPurple!9,minimum width=7.8cm,
    minimum height=0.42cm,inner sep=0pt,font=\footnotesize,text=ChalmersPurple]
    at (7.2,2.43) {P2P wrapper + bindings};

  \draw[<->,ChalmersPurple!75,line width=1pt] (7.2,2.17)--(7.2,1.37);
  \node[anchor=west,font=\footnotesize,text=muted]
    at (7.4,1.77) {Unix-domain sockets};

  % Nesting denotes a library inside the Go daemon, not another process.
  \fill[ChalmersCyan!7,rounded corners=3mm]
    (3.1,0.17) rectangle (11.3,1.32);
  \node[font=\small\bfseries,text=ChalmersPurple] at (7.2,1.05)
    {p2pd\quad\textbar\quad Go};
  \node[box,draw=none,fill=white,minimum width=7.4cm,minimum height=0.5cm,
    inner sep=0pt] at (7.2,0.57) {};
  \node[font=\small\bfseries,text=ChalmersPurple] at (4.8,0.57) {libp2p};
  \node[quiet,text=ink] at (8.3,0.57) {Connects peers securely.};
\end{visual}
\note{This is a map of components; the following slides explain their interactions.
The optimizer coordinates local training and averaging. The DHT stores discovery and
progress metadata, matchmaking forms temporary groups, and all-reduce computes a
weighted gradient average. The shared networking interface is Python's P2P wrapper
and daemon bindings. Python components can span several processes.
\par p2pd is a separate Go process containing libp2p. The single arrow represents local
RPCs and streams over Unix-domain sockets: Python connects to the daemon, and incoming
streams can reach a Python client listener. Lower-level bindings also support IPv4/TCP;
Hivemind creates Unix sockets by default. libp2p handles authenticated, encrypted peer
connections. Gradients bypass DHT storage.
\par Sources: p2p/p2p\_daemon.py:42,165--166,227,255;\\
p2p/p2p\_daemon\_bindings/control.py:57--64,157--165; setup.py:113--116;
progress\_tracker.py:220,247; optimizer.py:528,613--630;
key\_manager.py:62--80; averager.py:428--464,543--554.}
\end{frame}

\begin{frame}{Hivemind's training loop}
\begin{visual}
  \node[heading] at (2.3,5.15) {1\quad Train locally};
  \node[heading] at (7.3,5.15) {2\quad Average in a group};
  \node[heading] at (12.15,5.15) {3\quad Update locally};
  \draw[ChalmersPurple!45,fill=ChalmersPurple!3,rounded corners=3mm]
    (0.3,1.65) rectangle (4.3,4.55);
  \node[heading,text=ChalmersPurple] at (2.3,4.1) {Peer 1};
  \node[box,draw=muted!55,fill=white,inner sep=1.5mm,minimum width=12mm,
    minimum height=8mm] (data) at (1.25,3.0) {data};
  \node[box,fill=white,inner sep=1.5mm,minimum width=14mm,minimum height=10mm,
    font=\small] (model) at (3.15,3.0) {full\\model};
  \draw[flow] (data)--(model);
  \node[quiet] at (2.3,2.12) {compute a local gradient};

  \draw[ChalmersPurple!50,dashed,rounded corners=4mm]
    (5.15,1.65) rectangle (9.55,4.55);
  \node[peer,minimum size=11mm] (p1) at (6.15,3.8) {P1};
  \node[peer,minimum size=11mm] (p2) at (8.55,3.8) {P2};
  \node[peer,minimum size=11mm] (p3) at (7.35,2.4) {P3};
  \draw[dataflow] (p1)--(p2);
  \draw[dataflow] (p2)--(p3);
  \draw[dataflow] (p3)--(p1);
  \node[quiet,text=dat] at (7.35,1.23) {peers exchange gradients};
  \draw[flow] (4.4,3.1)--(5.05,3.1);

  \draw[ChalmersPurple!45,fill=ChalmersPurple!3,rounded corners=3mm]
    (10.35,1.65) rectangle (14.0,4.55);
  \node[heading,text=ChalmersPurple] at (12.175,4.1) {Peer 1};
  \node[box,draw=dat!65,fill=white,inner sep=2mm,minimum width=22mm,
    minimum height=10mm] (updated) at (12.175,3.0) {updated model};
  \draw[->,dat,line width=1.2pt] (9.7,3.0)--(updated.west);
  \node[quiet] at (12.175,2.12) {apply average gradient};

  \draw[flow,rounded corners=2mm] (12.175,1.55)--(12.175,0.65)--(2.3,0.65)--(2.3,1.55);
  \node[font=\footnotesize,text=ChalmersPurple,fill=white,inner sep=1mm] at (7.2,0.65) {repeat};
  \node[quiet] at (7.2,0.05) {Every peer keeps its own data and full model\quad\textbar\quad no central aggregation server};
\end{visual}
\note{Follow Peer 1 from left to right; it is the same machine in all three stages.
First, its own data and full model produce a local gradient, a vector describing a training
update. Second, it joins ready peers in a temporary group and exchanges gradients with them.
The circles P1,P2,P3 are Peers 1,2,3, all doing the same job. Third, Peer 1 applies the group
average to its local model and continues training. Its groupmates do the same. This is the
basic repeated operation Hivemind enables without a central aggregation server.
\par Unlike a server-based FL diagram, the middle is a group of training peers, not an extra
service. Later averaging operations assemble groups again, although membership need not change.
Different groups need not have identical model states. Raw data stays in the local training
pipeline; this does not guarantee gradient privacy. Bootstrap contacts help enter the network,
not coordinate learning. Sources: HIVEMIND\_EXPLAINED sections 1,4,6--8; HANDOFF 4.10.}
\end{frame}

% \begin{frame}{Hivemind's training loop}
% \begin{visual}
%   \node[heading] at (2.2,5.15) {Train at different speeds};
%   \node[heading] at (6.65,5.15) {One shared target};
%   \node[heading] at (11.55,5.15) {Average + update};

%   % Illustrative local sample contributions, not individual batch quotas.
%   \foreach \name/\yy/\width/\samples in {
%     P1/4.45/2.4/600,P2/3.86/1.6/400,P3/3.27/0.8/200} {
%     \node[font=\footnotesize\bfseries,text=ChalmersPurple,anchor=east]
%       at (0.83,\yy) {\name};
%     \fill[ChalmersCyan!65,rounded corners=0.7mm]
%       (1.0,\yy-0.13) rectangle (1.0+\width,\yy+0.13);
%     \node[font=\footnotesize,text=muted,anchor=west]
%       at (3.58,\yy) {\samples};
%   }
%   \draw[flow] (4.45,3.86)--(5.15,3.86);
%   \node[box,draw=none,fill=ChalmersPurple!6,minimum width=2.65cm,
%     minimum height=1.75cm,inner sep=0pt] at (6.65,3.86) {};
%   \node[font=\LARGE\bfseries,text=ChalmersPurple] at (6.65,4.09) {1,200};
%   \node[quiet] at (6.65,3.53) {samples / epoch};
%   \draw[flow] (8.05,3.86)--(8.8,3.86);

%   \draw[ChalmersPurple!40,dashed,rounded corners=3mm]
%     (8.97,2.99) rectangle (14.0,4.78);
%   \node[peer,minimum size=8mm] (t1) at (9.9,4.24) {P1};
%   \node[peer,minimum size=8mm] (t2) at (13.07,4.24) {P2};
%   \node[peer,minimum size=8mm] (t3) at (11.48,3.52) {P3};
%   \draw[dataflow,line width=0.8pt] (t1)--(t2);
%   \draw[dataflow,line width=0.8pt] (t2)--(t3);
%   \draw[dataflow,line width=0.8pt] (t3)--(t1);
%   \node[quiet] at (2.2,2.78) {local samples accumulated};
%   \node[quiet] at (11.55,2.78) {temporary averaging group};
%   \draw[flow] (12.8,2.39)--(1.6,2.39);
%   \node[font=\footnotesize,text=ChalmersPurple,fill=white,inner sep=1mm]
%     at (7.2,2.39) {Next epoch: train + form groups again};

%   % Optional redistribution is distinct from the normal epoch cycle above.
%   \node[font=\small\bfseries,text=ChalmersPurple] at (1.95,1.5) {Optional sharding};
%   \node[quiet] at (1.95,1.04) {off by default};
%   \draw[ChalmersPurple!40,dashed,rounded corners=2mm]
%     (3.65,0.65) rectangle (6.85,1.88);
%   \node[quiet,font=\scriptsize,text=muted,fill=white,inner sep=0.5mm]
%     at (5.25,1.88) {one group};
%   \foreach \xx/\member in {4.17/1,4.89/2,5.61/3,6.33/4}
%     \node[peer,minimum size=5.5mm,font=\scriptsize\bfseries] at (\xx,1.24) {\member};
%   \draw[flow] (7.0,1.24)--(7.85,1.24);
%   \foreach \xx/\bits/\member in {8.65/00/1,10.15/01/2,11.65/10/3,13.15/11/4} {
%     \node[box,draw=ChalmersPurple!25,fill=ChalmersPurple!3,minimum width=1.12cm,
%       minimum height=1.28cm,inner sep=0pt] at (\xx,1.24) {};
%     \node[font=\footnotesize\ttfamily,text=muted] at (\xx,1.61) {\bits};
%     \node[peer,minimum size=5.5mm,font=\scriptsize\bfseries] at (\xx,1.02) {\member};
%   }
%   \node[quiet] at (10.9,0.25) {different discovery buckets next round};
%   \node[quiet] at (3.7,0.25) {Default: groupmates may repeat};
% \end{visual}
% \note{An epoch is a shared training-sample target, not a data pass or a per-peer quota.
% Faster peers contribute more minibatches: 600+400+200=1200 is an illustration.
% Peers retain local data and full models. The progress tracker exchanges sample counts
% and speeds through the DHT, estimates epoch end, and schedules matchmaking near it;
% this is not an exact barrier. Temporary groups average gradients, peers update, then
% later averaging operations run matchmaking again. Their membership may repeat.
% \par The separate four-peer example shows optional sharding. Group assembly assigns
% members distinct next discovery buckets using the shared group ID and member positions.
% The example has enough suffix bits for four distinct assignments. These boxes are
% rendezvous buckets, not assembled groups. Ordinary co-members separate next round;
% an attacker can compute the same assignments and follow. Empty group bits make this
% a no-op by default. Next keys change on assembly, not at a global epoch barrier.
% \par Sources: optimizer.py:54--67,559--567; progress\_tracker.py:128--149,308--322;
% averager.py:120,427--459,501--512; key\_manager.py:94--104;
% matchmaking.py:387,404; HANDOFF 4.4,4.8.}
% \end{frame}

\begin{frame}{How a group forms}
\begin{visual}
  \node[heading,text=ChalmersPurple] at (1.55,5.2) {Peer 1};
  \node[box,draw=ctl!50,fill=ctl!6,inner sep=2mm,font=\small\bfseries]
    at (5.45,5.2) {DHT replicas};
  \node[heading,text=ChalmersPurple] at (9.1,5.2) {Peer 2};
  \node[heading,text=ChalmersPurple] at (12.85,5.2) {Peer 3};
  \foreach \xx in {1.55,5.45,9.1,12.85}
    \draw[muted!45,dashed] (\xx,4.8)--(\xx,0.65);
  \draw[->,muted!65] (0.3,4.45)--(0.3,1.1);
  \node[quiet,rotate=90] at (0.08,2.75) {time};
  \draw[control] (9.1,4.25)--node[above,font=\footnotesize]{availability + deadline}(5.45,4.25);
  \draw[control] (1.55,3.6)--node[above,font=\footnotesize]{lookup job / bucket}(5.45,3.6);
  \draw[control] (5.45,2.95)--node[above,font=\footnotesize]{candidates + deadlines}(1.55,2.95);
  \node[quiet,text=ctl,align=left,anchor=west] at (9.3,3.3) {earliest visible deadline\\becomes candidate leader};
  \filldraw[fill=ctl!12,draw=ctl!60] (8.96,0.7) rectangle (9.24,2.6);
  \draw[control] (1.55,2.25)--node[above,font=\footnotesize]{join request}(8.96,2.25);
  \draw[control] (12.85,1.75)--node[above,font=\footnotesize]{join request}(9.24,1.75);
  \draw[control] (8.96,1.2)--node[above,font=\footnotesize]{member list + group ID}(1.55,1.2);
  \draw[control] (9.24,0.85)--node[above,font=\footnotesize]{member list}(12.85,0.85);
  \node[box,draw=none,fill=ChalmersCream,minimum width=13.6cm,inner sep=2mm] at (7.2,0.15)
    {The temporary leader checks requests and admits members.};
\end{visual}
\note{Read this as a sequence diagram, downward in time. Only selected messages are shown;
all searching peers can publish advertisements. A job and optional bucket identify a DHT
rendezvous key. Records advertise a PeerID, availability and an expiry that also serves as
the intended matchmaking deadline. Peers look up candidates and prefer the earliest visible
deadline, with PeerID as tie-breaker. There is no consensus election or guaranteed common view.
\par Peer 2 is the candidate leader in this example. Others send it direct join RPCs. It checks
compatibility and protocol fields, admits followers, then distributes the group ID and member
list when full or when its deadline permits. The leader assembles the group; it is not the
sole gradient reducer. Groups form again for later averaging. Moving to different DHT buckets
is optional: the default empty group bits leave the rendezvous key unchanged. Sources:
key\_manager.py:43--104; matchmaking.py:262--405,469--498; HANDOFF sections 1,4.3,4.8.}
\end{frame}

\begin{frame}{All-reduce: averaging within the group}
\begin{visual}
  \node[heading] at (2.05,5.15) {1\quad Split};
  \node[heading] at (7.2,5.15) {2\quad Average each slice};
  \node[heading] at (12.35,5.15) {3\quad Reassemble};
  \foreach \peer/\yy/\a/\b/\c in {1/4.05/2/4/6,2/2.95/4/6/8,3/1.85/6/8/10} {
    \node[quiet,anchor=east] at (0.95,\yy) {Peer \peer};
    \foreach \xx/\val/\col in {1.45/\a/ChalmersCyan,2.25/\b/ChalmersGreen,3.05/\c/ChalmersPurple} {
      \node[draw=\col!65,fill=\col!12,minimum width=7.5mm,minimum height=7mm,
        inner sep=0pt,font=\normalsize\bfseries] at (\xx,\yy) {\val};
    }
    \node[quiet,anchor=east] at (11.3,\yy) {Peer \peer};
    \foreach \xx/\val/\col in {11.8/4/ChalmersCyan,12.6/6/ChalmersGreen,13.4/8/ChalmersPurple} {
      \node[draw=\col!65,fill=\col!12,minimum width=7.5mm,minimum height=7mm,
        inner sep=0pt,font=\normalsize\bfseries] at (\xx,\yy) {\val};
    }
  }
  \foreach \peer/\yy/\vals/\mean/\col in {
    1/4.05/{2,4,6}/4/ChalmersCyan,
    2/2.95/{4,6,8}/6/ChalmersGreen,
    3/1.85/{6,8,10}/8/ChalmersPurple} {
    \node[box,draw=\col!65,fill=\col!8,minimum width=4.0cm,minimum height=8mm,
      inner sep=2mm] at (7.2,\yy)
      {\textbf{Peer \peer}\quad $\operatorname{mean}(\vals)=\mean$};
  }
  \draw[dataflow,->] (3.65,2.95)--(4.85,2.95);
  \draw[dataflow,->] (9.45,2.95)--(10.1,2.95);
  \node[quiet] at (2.05,0.95) {local gradients};
  \node[quiet] at (7.2,0.95) {a different reducer for each slice};
  \node[quiet] at (12.45,0.95) {same group average};
  \node[quiet] at (7.2,0.15)
    {All-reduce: every group member obtains the aggregate\quad\textbar\quad equal-weight illustration};
\end{visual}
\note{All-reduce means all group members obtain the aggregate. The toy vectors make the
operation concrete: each peer contributes three numbers, and the group average is [4,6,8].
Peer 1 reduces the blue slice, Peer 2 the green slice, Peer 3 the purple slice. Each sender
contacts the assigned reducer directly; this is not ring forwarding. The same group carries
out all three stages. The DHT stores metadata; peers exchange gradient slices directly using
libp2p, which may use transport relays underneath.
\par Real gradients are long vectors, slices can have unequal sizes, and Hivemind uses a
weighted mean. Reducers return averaged-minus-local deltas; peers reconstruct the result and
the optimizer applies the update. Equal slices and equal weights here are explanatory choices.
The temporary group leader is not a central aggregation server. Sources: allreduce.py:93--111,
164--216,335--356; partition.py:242--260; HANDOFF 4.10--4.11.}
\end{frame}

\begin{frame}{Two IDs}
\begin{visual}
  \node[heading] at (3.5,5.12) {Place records};
  \node[heading] at (10.9,5.12) {Contact peers};
  \node[box,draw=ctl!40,fill=ctl!5,minimum width=5.8cm,minimum height=2.75cm]
    at (3.5,3.3) {};
  \node[box,minimum width=5.8cm,minimum height=2.75cm] at (10.9,3.3) {};
  \node[heading,text=ctl] at (3.5,4.28) {Hivemind DHT: node\_id};
  \node[heading,text=ChalmersPurple] at (10.9,4.28) {libp2p: PeerID};

  \node[labeltext] at (3.5,3.5) {160-bit position in the DHT};
  \node[quiet] at (3.5,2.85) {By default, the 5 nearest nodes\\(XOR distance) store each record.};
  \node[labeltext,text=ctl,font=\small\bfseries] at (3.5,2.25) {chosen by the peer};
  \node[labeltext] at (10.9,3.5) {identity derived from a public key};
  \node[quiet] at (10.9,2.85) {libp2p resolves its address\\and authenticates the connection};
  \node[labeltext,text=ChalmersPurple,font=\small\bfseries] at (10.9,2.25) {key ownership checked};

  % Both identifiers belong to the same peer; there is no ID-conversion arrow.
  \draw[muted!35,rounded corners=2mm,fill=white] (4.0,0.43) rectangle (10.05,1.78);
  \node[font=\footnotesize\bfseries,text=muted] at (7.025,1.5) {One peer, two IDs};
  \node[box,draw=muted!40,fill=white,inner sep=2mm,minimum width=2.2cm]
    (key) at (1.6,0.88) {record key};
  \node[box,draw=ctl!55,fill=ctl!5,inner sep=2mm,minimum width=2.9cm]
    (pos) at (5.62,0.88) {nearest \texttt{node\_id}};
  \node[box,inner sep=2mm,minimum width=2.1cm]
    (endpoint) at (8.75,0.88) {PeerID};
  \node[box,inner sep=2mm,minimum width=2.5cm]
    (rpc) at (12.5,0.88) {connect + RPC};
  \draw[control] (key)--(pos);
  \draw[flow] (endpoint)--(rpc);
  \node[quiet] at (7.2,0.02)
    {The routing table stores both IDs; libp2p connects using PeerID.};
\end{visual}
\note{The DHT uses node\_id as a 160-bit XOR coordinate for routing and replica placement.
It normally stores each record on the five nearest nodes, trying other candidates if stores
fail. Additional cached copies can exist.
\par PeerID identifies a libp2p endpoint through its public key. libp2p resolves network
addresses and authenticates key ownership; PeerID itself is not an IP address.
Hivemind builds its placement overlay on this networking layer, so the IDs have different jobs.
The routing table stores both for each peer: select a node\_id, then use its associated
PeerID to connect. One ID is not computed from the other. The diagram shows one holder;
real lookups and replication involve several nodes.
\par Separate roles do not require unrelated IDs. Here node\_id is random by default but
can be chosen independently, without a key-derived binding to PeerID. Transport authentication
therefore does not validate DHT placement or provide Sybil resistance.
Sources: dht/routing.py:34--64,111; dht/node.py:105,171,411--462;
dht/protocol.py:149--158; p2p\_daemon\_bindings/control.py:370.}
\end{frame}
